% =============================================================================
\chapter{Fundamentação Teórica Inicial}
% =============================================================================

\section{Aprendizado por reforço}

Aprendizado por reforço é uma área de aprendizado de máquina voltada à tomada
de decisão sequencial.
Um agente interage com um ambiente, observa estados, executa ações e recebe
recompensas.
O objetivo é aprender uma política que maximize algum critério de retorno esperado
ao longo do tempo.

Um problema de aprendizado por reforço pode ser descrito pela interação entre
agente e ambiente:

\begin{equation}
s_t \rightarrow a_t \rightarrow r_{t+1} \rightarrow s_{t+1}.
\end{equation}

Nessa interação, $s_t$ é o estado observado no instante $t$, $a_t$ é a ação
escolhida pelo agente, $r_{t+1}$ é a recompensa recebida após a execução da
ação e $s_{t+1}$ é o próximo estado do ambiente.
O agente busca aprender como escolher ações que maximizem recompensas futuras,
considerando que uma decisão tomada no presente pode afetar estados e recompensas
posteriores.

Essa formulação é relevante para a pesquisa porque o projeto investiga agentes
autônomos capazes de tomar decisões em ambientes simulados.
Antes de integrar aprendizado diretamente ao agente, é necessário definir
claramente como estados, ações, recompensas, objetivos e transições serão
representados.
Por isso, a base teórica em Processos de Decisão de Markov, funções de valor,
diferença temporal e aprendizado por reforço multiobjetivo é importante para
as etapas futuras do trabalho.

\section{Processos de Decisão de Markov}

Um Processo de Decisão de Markov, ou MDP, é uma formulação matemática para
problemas de tomada de decisão sequencial sob incerteza.
Essa formulação fornece a base formal para grande parte dos algoritmos de
aprendizado por reforço.

Um MDP pode ser definido como uma tupla:

\begin{equation}
M = \langle S, A, P, R, \gamma \rangle.
\end{equation}

Nessa definição:

\begin{itemize}
\item $S$ é o conjunto de estados possíveis do ambiente;
\item $A$ é o conjunto de ações disponíveis para o agente;
\item $P(s' \mid s,a)$ é a função de transição, que define a probabilidade de ir para o estado $s'$ após executar a ação $a$ no estado $s$;
\item $R(s,a,s')$ é a função de recompensa, que define o retorno imediato associado à transição de $s$ para $s'$ por meio da ação $a$;
\item $\gamma$ é o fator de desconto, com $0 \leq \gamma \leq 1$, usado para controlar a importância de recompensas futuras.
\end{itemize}

A propriedade de Markov estabelece que o próximo estado depende apenas do
estado atual e da ação executada, e não de todo o histórico anterior.
Em outras palavras, se o estado atual contém toda a informação necessária para
prever a dinâmica futura do ambiente, então o histórico completo pode ser ignorado
para fins de decisão.

Formalmente, essa propriedade pode ser escrita como:

\begin{equation}
P(s_{t+1} \mid s_t,a_t,s_{t-1},a_{t-1},\dots,s_0,a_0)
=
P(s_{t+1} \mid s_t,a_t).
\end{equation}

Em um ambiente discreto inicial, como um GridWorld, essa propriedade significa
que o estado precisa conter todas as informações relevantes para a decisão do
agente.
Por exemplo, a posição do agente, a posição de recursos, a posição de objetivos,
o conjunto de obstáculos e variáveis booleanas relevantes devem ser suficientes
para determinar a dinâmica futura do ambiente.

Uma política define o comportamento do agente.
Em uma política estocástica, $\pi(a \mid s)$ representa a probabilidade de
escolher a ação $a$ quando o agente está no estado $s$:

\begin{equation}
\pi(a \mid s).
\end{equation}

Em uma política determinística, a política pode ser representada simplesmente
como uma função que mapeia estados para ações:

\begin{equation}
\pi(s) = a.
\end{equation}

O objetivo do agente é encontrar uma política que maximize o retorno esperado.
O retorno descontado a partir do instante $t$ é definido como:

\begin{equation}
G_t =
R_{t+1}
+
\gamma R_{t+2}
+
\gamma^2 R_{t+3}
+
\dots.
\end{equation}

De forma compacta:

\begin{equation}
G_t =
\sum_{k=0}^{\infty} \gamma^k R_{t+k+1}.
\end{equation}

O fator de desconto $\gamma$ controla a importância de recompensas futuras.
Quando $\gamma$ é próximo de zero, o agente prioriza recompensas imediatas.
Quando $\gamma$ é próximo de um, o agente considera recompensas de longo
prazo.

A função de valor de estado mede o retorno esperado ao seguir uma política
$\pi$ a partir de um estado $s$:

\begin{equation}
V_{\pi}(s) =
E_{\pi}[G_t \mid s_t = s].
\end{equation}

A função de valor de ação mede o retorno esperado ao executar uma ação $a$ em
um estado $s$ e, a partir daí, seguir a política $\pi$:

\begin{equation}
Q_{\pi}(s,a) =
E_{\pi}[G_t \mid s_t = s, a_t = a].
\end{equation}

Para este projeto, os MDPs são relevantes porque fornecem a formulação
matemática da interação entre agente e ambiente.
Eles permitem descrever a dinâmica sequencial em termos de estados, ações,
transições e recompensas.
Em etapas futuras, essa formulação permitirá comparar agentes simbólicos com
agentes baseados em aprendizado.

\section{Q-Learning como base conceitual}

Q-Learning é um algoritmo clássico de aprendizado por reforço usado para
aprender uma função de valor de ação.
Embora não seja o método principal pretendido para esta pesquisa, ele é relevante
como fundamento conceitual, pois introduz ideias centrais que também aparecem em
métodos mais avançados: função de valor, equação de Bellman, diferença temporal e
atualização baseada em experiência.

A função ótima de valor de ação, denotada por $Q^*(s,a)$, representa o retorno
esperado ao executar a ação $a$ no estado $s$ e seguir uma política ótima a
partir do próximo estado.
Essa função satisfaz a equação de otimalidade de Bellman:

\begin{equation}
Q^*(s,a) =
\sum_{s'} P(s' \mid s,a)
\left[
R(s,a,s')
+
\gamma \max_{a'} Q^*(s',a')
\right].
\end{equation}

Essa equação expressa que o valor ótimo de uma ação depende de duas partes:
a recompensa imediata obtida após sua execução e o melhor valor futuro
possível a partir do próximo estado.

Na prática, o agente geralmente não conhece a função de transição
$P(s' \mid s,a)$ nem a função de recompensa completa do ambiente.
Por isso, Q-Learning estima os valores de ação a partir de amostras de experiência.
Uma experiência pode ser representada por:

\begin{equation}
(s_t, a_t, r_{t+1}, s_{t+1}).
\end{equation}

A regra de atualização do Q-Learning é:

\begin{equation}
Q(s_t,a_t)
\leftarrow
Q(s_t,a_t)
+
\alpha
\left[
r_{t+1}
+
\gamma \max_{a'} Q(s_{t+1},a')
-
Q(s_t,a_t)
\right].
\end{equation}

Nessa equação, $\alpha$ é a taxa de aprendizado, $\gamma$ é o fator de
desconto, $r_{t+1}$ é a recompensa observada e
$\max_{a'} Q(s_{t+1},a')$ representa a melhor estimativa de valor futuro no
estado seguinte.

O termo entre colchetes pode ser interpretado como erro de diferença temporal:

\begin{equation}
\delta_t =
r_{t+1}
+
\gamma \max_{a'} Q(s_{t+1},a')
-
Q(s_t,a_t).
\end{equation}

Assim, a atualização também pode ser escrita como:

\begin{equation}
Q(s_t,a_t)
\leftarrow
Q(s_t,a_t)
+
\alpha \delta_t.
\end{equation}

A intuição é que o agente ajusta sua estimativa atual na direção de um alvo
temporal.
Esse alvo é composto pela recompensa imediata mais a melhor estimativa de
valor futuro:

\begin{equation}
\text{target} =
r_{t+1}
+
\gamma \max_{a'} Q(s_{t+1},a').
\end{equation}

Q-Learning é considerado um método \textit{off-policy}.
Isso significa que ele pode aprender o valor da política ótima mesmo
quando as ações são escolhidas por uma política exploratória.
Por exemplo, o agente pode escolher ações usando uma política
$\varepsilon$-\textit{greedy}, mas atualizar sua função $Q$ usando o melhor
valor futuro possível, dado por $\max_{a'} Q(s_{t+1},a')$.

Para este plano, Q-Learning não deve ser tratado como técnica
principal da pesquisa nem como componente previsto da arquitetura
final.
Sua função é servir como base conceitual para entender como
funções de valor podem ser estimadas a partir da interação
com o ambiente.

A técnica de interesse para a continuidade do projeto é aprendizado por
reforço multiobjetivo.
Nesse caso, em vez de aprender uma única função de valor escalar, o
agente deve lidar com múltiplos critérios de decisão, representados por
diferentes sinais de recompensa ou diferentes funções de valor.
Assim, Q-Learning aparece neste plano apenas como ponto de partida
conceitual para compreender funções de valor e atualização temporal,
enquanto MORL representa a direção metodológica pretendida para a camada
de aprendizado.

\section{Multi-armed bandits}

Uma base importante estudada no início da pesquisa foi o problema dos
\textit{multi-armed bandits}.
Nesse problema, o agente escolhe repetidamente entre $k$ ações, cada uma
com uma recompensa de média desconhecida.
O valor real de uma ação é:

\begin{equation}
q_*(a) = E[R_t \mid A_t = a].
\end{equation}

Como esse valor é desconhecido, o agente mantém uma estimativa $Q_t(a)$.
A estimativa por média amostral é:

\begin{equation}
Q_t(a) =
\frac{\sum_{i=1}^{t-1} R_i I(A_i = a)}
{\sum_{i=1}^{t-1} I(A_i = a)}.
\end{equation}

Essa formulação é importante porque introduz conceitos fundamentais como
estimativa de valor, exploração, aproveitamento e atualização incremental em
um cenário mais simples do que MDPs completos.
Enquanto bandits não modelam transições explícitas entre estados, MDPs
incorporam o efeito das ações sobre estados futuros.
Por isso, bandits são uma porta de entrada conceitual para o
estudo de aprendizado por reforço.

\section{Atualização incremental}

A atualização incremental evita recalcular a média completa a cada nova
observação.
A regra geral é:

\begin{equation}
Q_{n+1}(a) = Q_n(a) + \alpha_n(R_n - Q_n(a)).
\end{equation}

O termo $R_n - Q_n(a)$ representa o erro de previsão.
Quando $\alpha_n = 1/n$, a atualização equivale à média amostral.
Quando $\alpha$ é constante, o método dá maior peso a recompensas recentes,
o que é útil em ambientes não estacionários.

Essa ideia reaparece no Q-Learning.
A diferença é que, no caso de Q-Learning, o erro de atualização inclui
não apenas a recompensa imediata, mas também uma estimativa do melhor
valor futuro possível.

\section{Exploração e aproveitamento}

O dilema exploração versus aproveitamento é central em aprendizado por
reforço.
Explorar significa testar ações pouco conhecidas.
Aproveitar significa escolher a ação com melhor valor estimado.

Uma política simples para equilibrar os dois comportamentos é
$\varepsilon$-\textit{greedy}.
Com probabilidade $1-\varepsilon$, o agente escolhe a melhor ação conhecida.
Com probabilidade $\varepsilon$, escolhe uma ação aleatória.

Outra estratégia estudada foi UCB, que escolhe ações considerando tanto o
valor estimado quanto a incerteza:

\begin{equation}
A_t =
\arg\max_a
\left[
Q_t(a) + c\sqrt{\frac{\ln t}{N_t(a)}}
\right].
\end{equation}

Aqui, $N_t(a)$ indica quantas vezes a ação $a$ foi escolhida até o instante
$t$, e $c$ controla a intensidade da exploração.

No projeto, exploração e aproveitamento serão relevantes principalmente nas
etapas futuras de aprendizado.
A camada simbólica atual não aprende por tentativa e erro; ela usa
decomposição explícita de objetivos e ações.
Entretanto, quando métodos de aprendizado forem introduzidos, será necessário
definir como o agente explora ações sem violar restrições simbólicas
importantes.

\section{Planejamento simbólico}

Planejamento simbólico representa estados, tarefas, ações e restrições por
estruturas discretas e interpretáveis.
Em vez de aprender diretamente todo o comportamento, o agente raciocina sobre
pré-condições, efeitos e relações de decomposição.
Essa representação permite incorporar conhecimento de domínio e impedir
decisões incompatíveis com regras narrativas ou de gameplay.

No presente trabalho, as restrições simbólicas são tratadas como
\textit{hard constraints}.
Preferências como segurança, rapidez ou economia de recursos são tratadas
separadamente como critérios flexíveis.
Essa separação é essencial: uma preferência pode mudar em tempo de execução,
enquanto uma regra como não utilizar um recurso indisponível deve permanecer
obrigatória.

\section{Planejamento HTN}

HTN, ou \textit{Hierarchical Task Network Planning}, decompõe tarefas
compostas em subtarefas até produzir tarefas primitivas executáveis
\cite{erol1994,georgievski2014,nau2003shop2}.
Um domínio HTN define tarefas, métodos de decomposição, pré-condições
e restrições.
Para uma tarefa composta $\tau_t$ e estado $s_t$, o conjunto de métodos
aplicáveis é:

\begin{equation}
M_{\mathrm{valid}}(s_t,\tau_t)
=
\left\{m \in M(\tau_t) \mid \mathrm{pre}(m,s_t)=\mathrm{true}\right\}.
\end{equation}

Na arquitetura proposta, o HTN não escolhe imediatamente o primeiro método
aplicável.
Ele produz o conjunto de alternativas válidas e o entrega ao componente
aprendido.
Assim, o modelo simbólico mantém a responsabilidade pela validade, enquanto
a preferência entre alternativas válidas pode ser adaptada por experiência.

\section{Codificação de preferências}

GOAP representa objetivos, ações, pré-condições, efeitos e custos para buscar
sequências que transformem o estado corrente em um estado desejado
\cite{orkin2005,orkin2006}.
Neste trabalho, GOAP não será introduzido como um segundo planejador entre
HTN e MORL. Suas relações entre objetivos, condições, efeitos e recursos
podem, contudo, inspirar a representação de contexto usada por um codificador
de preferências.

O codificador transforma objetivo, estado, perfil e contexto em um vetor
normalizado de preferências:

\begin{equation}
\mathbf{w}_t = E(g_t,s_t,p,c_t),
\qquad
w_{i,t}\geq 0,
\qquad
\sum_i w_{i,t}=1.
\end{equation}

Nessa expressão, $E$ é o codificador, $g_t$ é o objetivo corrente, $s_t$ é o
estado, $p$ representa um perfil comportamental, $c_t$ é o contexto e $w_{i,t}$
é o peso do objetivo $i$ no instante $t$.
Perfis fixos e regras explícitas constituem baselines reproduzíveis.
A principal alternativa estrutural é um grafo relacional inspirado em GOAP,
sem executar sua busca de planejamento; um modelo de aprendizado profundo será
avaliado apenas se os dados, o orçamento e os resultados dos baselines justificarem
essa extensão.
Assim, a comparação entre codificadores é uma parte auxiliar da pesquisa, enquanto
a contribuição central permanece a seleção MORL de métodos HTN válidos.

\section{Aprendizado por reforço multiobjetivo}

No MORL, a recompensa é um vetor
\cite{roijers2013,hayes2022morl}:

\begin{equation}
\mathbf{r}_t =
\left[r_{1,t},r_{2,t},\dots,r_{n,t}\right].
\end{equation}

Cada componente $r_{i,t}$ representa um objetivo, como sobrevivência,
progresso da missão, proteção de aliados, conservação de recursos ou tempo.
A política pode ser condicionada por $\mathbf{w}_t$, permitindo que uma única
aproximação represente diferentes compromissos entre objetivos.
Trabalhos sobre pesos dinâmicos e transferência de políticas fornecem
referências para essa formulação \cite{abels2019,barreto2017successor,alegre2022ols,alegre2026gpi}.

Para cada método válido, a política estima um vetor de valor:

\begin{equation}
\mathbf{Q}(s_t,\tau_t,m,\mathbf{w}_t)
=
\left[Q_1,Q_2,\dots,Q_n\right].
\end{equation}

Aqui, $\tau_t$ é a tarefa composta corrente, $m$ é um método candidato e cada
$Q_i$ estima o retorno descontado do objetivo $i$ ao selecionar $m$. A
dependência explícita de $\mathbf{w}_t$ permite à aproximação aprender valores
ou políticas condicionadas pela preferência, sem transferir ao HTN a decisão de
compromisso entre objetivos.

Sob escalarização linear, a seleção é:

\begin{equation}
 m_t^* =
 \arg\max_{m\in M_{\mathrm{valid}}(s_t,\tau_t)}
 \mathbf{w}_t^{\top}
 \mathbf{Q}(s_t,\tau_t,m,\mathbf{w}_t).
\end{equation}

A máscara de métodos válidos restringe exploração e inferência às alternativas
permitidas pelo HTN. Desse modo, o MORL aprende preferências entre decisões
semanticamente aceitáveis, em vez de aprender diretamente ações de movimento,
mira ou animação.

\section{Decisão hierárquica e abstração temporal}

A seleção de um método HTN pode produzir uma sequência de subtarefas cuja
execução dura vários ticks.
Consequentemente, a decisão aprendida ocorre em uma escala temporal superior à
de ações primitivas e pode ser modelada de forma compatível com abstração
temporal e formulações semi-Markovianas \cite{sutton1999options,lin2024hipo}.
O retorno associado a uma escolha de método deve considerar os resultados obtidos
durante a execução da decomposição, até a próxima decisão hierárquica ou até uma
condição de interrupção.
Se o método selecionado dura $k$ passos primitivos, a unidade de atribuição de
crédito é:

\begin{equation}
\mathbf{R}^{\mathrm{method}}_t =
\sum_{i=0}^{k-1}\gamma^i\mathbf{r}_{t+i}.
\end{equation}

Nessa expressão, $\mathbf{r}_{t+i}$ é a recompensa vetorial no passo primitivo
$t+i$, $\gamma$ é o fator de desconto e $k$ é a duração da execução até a
próxima escolha de método, término ou interrupção.

\section{Arquitetura híbrida proposta}

A integração é definida pelo seguinte fluxo:

\begin{align}
(s_t,\tau_t)
&\longrightarrow \text{filtragem HTN}
\longrightarrow M_{\mathrm{valid}}(s_t,\tau_t),\\
(g_t,s_t,p,c_t)
&\longrightarrow E
\longrightarrow \mathbf{w}_t,\\
(M_{\mathrm{valid}},s_t,\tau_t,\mathbf{w}_t)
&\longrightarrow \text{MORL}
\longrightarrow m_t^*.
\end{align}

O método selecionado é decomposto em subtarefas, e as tarefas primitivas são
executadas por controladores tradicionais.
Após a execução, o ambiente produz um novo estado e um vetor de recompensas,
permitindo atualizar ou avaliar a política.
Essa divisão torna explícitas quatro responsabilidades: validade semântica,
preferência comportamental, utilidade aprendida e execução de baixo nível.

\begin{figure}[htbp]
\centering
\resizebox{0.94\textwidth}{!}{%
\begin{tikzpicture}[
    font=\small,
    >=Latex,
    node distance=8mm and 16mm,
    block/.style={
        rectangle, rounded corners=2mm, draw, thick, align=center,
        minimum height=10mm, minimum width=33mm
    },
    game/.style={block, draw=gamegreen, fill=gamegreen!8},
    htn/.style={block, draw=htnblue, fill=htnblue!8},
    morl/.style={block, draw=morlred, fill=morlred!8},
    runtime/.style={block, draw=envgray, fill=envgray!7},
    arrow/.style={->, thick}
]
    \node[game] (director) {Game / AI Director};
    \node[game, below left=10mm and 19mm of director] (task)
        {Main Task\\$T_{\mathrm{main}} = C_0$};
    \node[game, below right=10mm and 19mm of director] (weight)
        {Preference Vector\\$\mathbf{w}_t=[w_1,\ldots,w_d]$};
    \node[runtime, below=24mm of director] (state)
        {World State\\$WS_t$};
    \node[runtime, left=45mm of state] (sensors)
        {Environment / Sensors};

    \draw[arrow, gamegreen] (director) -- (task);
    \draw[arrow, gamegreen] (director) -- (weight);
    \draw[arrow, envgray, dashed] (sensors) -- (state);

    \node[htn, below=11mm of state] (compound)
        {Current Compound Task\\$C_t$};
    \draw[arrow, htnblue] (task) |- (compound);
    \draw[arrow, htnblue] (state) -- (compound);

    \node[htn, below=10mm of compound] (filter)
        {HTN Symbolic Filter\\Preconditions / Constraints};
    \node[htn, below=9mm of filter] (mask)
        {Feasible-Method Mask\\$\mathcal{M}_{\mathrm{valid}}(C_t,WS_t)$};
    \draw[arrow, htnblue] (compound) -- (filter);
    \draw[arrow, htnblue] (filter) -- (mask);

    \node[morl, below=10mm of mask] (morl)
        {Preference-Conditioned MORL\\$\mathbf{Q}_{\theta}(C_t,WS_t,M_i,\mathbf{w}_t)$};
    \draw[arrow, morlred] (mask) -- (morl);
    \draw[arrow, morlred, dashed] (weight.east)
        to[out=0,in=15] node[pos=.58,right] {preferences} (morl.east);

    \node[morl, below=10mm of morl] (selected)
        {Selected Method\\$M_t^*$};
    \draw[arrow, morlred] (morl) --
        node[right] {$\arg\max U_{\mathbf{w}_t}$} (selected);

    \node[htn, below=10mm of selected] (decompose)
        {HTN Decomposition\\Decompose only $M_t^*$};
    \node[htn, below=10mm of decompose] (network)
        {Resulting Task Network};
    \draw[arrow, htnblue] (selected) -- (decompose);
    \draw[arrow, htnblue] (decompose) -- (network);

    \node[htn, below left=11mm and 19mm of network] (primitive)
        {Primitive Tasks};
    \node[htn, below right=11mm and 19mm of network] (subcompound)
        {Nested Compound Task};
    \draw[arrow, htnblue] (network) -- (primitive);
    \draw[arrow, htnblue] (network) -- (subcompound);
    \draw[arrow, htnblue, dashed, rounded corners] (subcompound.east)
        -- ++(18mm,0) |- node[pos=.25,right,align=left]
        {next HTN\\decision point} (compound.east);

    \node[runtime, below=11mm of primitive] (plan)
        {Final Primitive Plan};
    \draw[arrow, envgray] (primitive) -- (plan);
\end{tikzpicture}%
}
\caption{Arquitetura proposta para seleção online de métodos HTN guiada por MORL.
O HTN determina a máscara de métodos aplicáveis; o MORL seleciona diretamente um
único método $M_t^*$ condicionado por $WS_t$ e $\mathbf{w}_t$. Apenas o método
selecionado é decomposto, e a seleção se repete em subtarefas compostas.}
\label{fig:htn-morl-architecture}
\end{figure}

